% Options for packages loaded elsewhere
\PassOptionsToPackage{unicode}{hyperref}
\PassOptionsToPackage{hyphens}{url}
%
\documentclass[
  10pt,
]{article}
\usepackage{amsmath,amssymb}
\usepackage{iftex}
\ifPDFTeX
  \usepackage[T1]{fontenc}
  \usepackage[utf8]{inputenc}
  \usepackage{textcomp} % provide euro and other symbols
\else % if luatex or xetex
  \usepackage{unicode-math} % this also loads fontspec
  \defaultfontfeatures{Scale=MatchLowercase}
  \defaultfontfeatures[\rmfamily]{Ligatures=TeX,Scale=1}
\fi
\usepackage{lmodern}
\ifPDFTeX\else
  % xetex/luatex font selection
  % Native editor uses its default main font.
  % Native editor uses its default sans font.
  % Native editor uses its default mono font.
\fi
% Use upquote if available, for straight quotes in verbatim environments
\IfFileExists{upquote.sty}{\usepackage{upquote}}{}
\IfFileExists{microtype.sty}{% use microtype if available
  \usepackage[]{microtype}
  \UseMicrotypeSet[protrusion]{basicmath} % disable protrusion for tt fonts
}{}
\makeatletter
\@ifundefined{KOMAClassName}{% if non-KOMA class
  \IfFileExists{parskip.sty}{%
    \usepackage{parskip}
  }{% else
    \setlength{\parindent}{0pt}
    \setlength{\parskip}{6pt plus 2pt minus 1pt}}
}{% if KOMA class
  \KOMAoptions{parskip=half}}
\makeatother
\usepackage{xcolor}
\usepackage[margin=0.8in]{geometry}
\usepackage{color}
\usepackage{fancyvrb}
\newcommand{\VerbBar}{|}
\newcommand{\VERB}{\Verb[commandchars=\\\{\}]}
\DefineVerbatimEnvironment{Highlighting}{Verbatim}{commandchars=\\\{\}}
% Add ',fontsize=\small' for more characters per line
\newenvironment{Shaded}{}{}
\newcommand{\AlertTok}[1]{\textcolor[rgb]{1.00,0.00,0.00}{\textbf{#1}}}
\newcommand{\AnnotationTok}[1]{\textcolor[rgb]{0.38,0.63,0.69}{\textbf{\textit{#1}}}}
\newcommand{\AttributeTok}[1]{\textcolor[rgb]{0.49,0.56,0.16}{#1}}
\newcommand{\BaseNTok}[1]{\textcolor[rgb]{0.25,0.63,0.44}{#1}}
\newcommand{\BuiltInTok}[1]{\textcolor[rgb]{0.00,0.50,0.00}{#1}}
\newcommand{\CharTok}[1]{\textcolor[rgb]{0.25,0.44,0.63}{#1}}
\newcommand{\CommentTok}[1]{\textcolor[rgb]{0.38,0.63,0.69}{\textit{#1}}}
\newcommand{\CommentVarTok}[1]{\textcolor[rgb]{0.38,0.63,0.69}{\textbf{\textit{#1}}}}
\newcommand{\ConstantTok}[1]{\textcolor[rgb]{0.53,0.00,0.00}{#1}}
\newcommand{\ControlFlowTok}[1]{\textcolor[rgb]{0.00,0.44,0.13}{\textbf{#1}}}
\newcommand{\DataTypeTok}[1]{\textcolor[rgb]{0.56,0.13,0.00}{#1}}
\newcommand{\DecValTok}[1]{\textcolor[rgb]{0.25,0.63,0.44}{#1}}
\newcommand{\DocumentationTok}[1]{\textcolor[rgb]{0.73,0.13,0.13}{\textit{#1}}}
\newcommand{\ErrorTok}[1]{\textcolor[rgb]{1.00,0.00,0.00}{\textbf{#1}}}
\newcommand{\ExtensionTok}[1]{#1}
\newcommand{\FloatTok}[1]{\textcolor[rgb]{0.25,0.63,0.44}{#1}}
\newcommand{\FunctionTok}[1]{\textcolor[rgb]{0.02,0.16,0.49}{#1}}
\newcommand{\ImportTok}[1]{\textcolor[rgb]{0.00,0.50,0.00}{\textbf{#1}}}
\newcommand{\InformationTok}[1]{\textcolor[rgb]{0.38,0.63,0.69}{\textbf{\textit{#1}}}}
\newcommand{\KeywordTok}[1]{\textcolor[rgb]{0.00,0.44,0.13}{\textbf{#1}}}
\newcommand{\NormalTok}[1]{#1}
\newcommand{\OperatorTok}[1]{\textcolor[rgb]{0.40,0.40,0.40}{#1}}
\newcommand{\OtherTok}[1]{\textcolor[rgb]{0.00,0.44,0.13}{#1}}
\newcommand{\PreprocessorTok}[1]{\textcolor[rgb]{0.74,0.48,0.00}{#1}}
\newcommand{\RegionMarkerTok}[1]{#1}
\newcommand{\SpecialCharTok}[1]{\textcolor[rgb]{0.25,0.44,0.63}{#1}}
\newcommand{\SpecialStringTok}[1]{\textcolor[rgb]{0.73,0.40,0.53}{#1}}
\newcommand{\StringTok}[1]{\textcolor[rgb]{0.25,0.44,0.63}{#1}}
\newcommand{\VariableTok}[1]{\textcolor[rgb]{0.10,0.09,0.49}{#1}}
\newcommand{\VerbatimStringTok}[1]{\textcolor[rgb]{0.25,0.44,0.63}{#1}}
\newcommand{\WarningTok}[1]{\textcolor[rgb]{0.38,0.63,0.69}{\textbf{\textit{#1}}}}
\setlength{\emergencystretch}{3em} % prevent overfull lines
\providecommand{\tightlist}{%
  \setlength{\itemsep}{0pt}\setlength{\parskip}{0pt}}
\setcounter{secnumdepth}{-\maxdimen} % remove section numbering
\usepackage{fvextra}
\fvset{breaklines=true,breakanywhere=true,fontsize=\scriptsize}
\setlength{\emergencystretch}{3em}
\ifLuaTeX
  \usepackage{selnolig}  % disable illegal ligatures
\fi
\usepackage{bookmark}
\IfFileExists{xurl.sty}{\usepackage{xurl}}{} % add URL line breaks if available
\urlstyle{same}
\hypersetup{
  hidelinks,
  pdfcreator={LaTeX via pandoc}}

\author{}
\date{}

\begin{document}

\section{Completed algebraic transport of ordinary zero-Hessian currents
in every pure
dimension}\label{completed-algebraic-transport-of-ordinary-zero-hessian-currents-in-every-pure-dimension}

Reviewed partial theorem, 4 October 2026. This is partial progress on
the original P06.2 recognition problem. Independent same-model review,
source integrity, finite diagnostics and compilation are not formal
verification or priority certification. No previous theorem, catalogue assertion, review verdict,
or hash is a premise.

\subsection{1. Statements and
conventions}\label{statements-and-conventions}

Let \(K\) be a complete nontrivially valued non-Archimedean field,
let \(M/K\) be an algebraic valued extension, and let \(L\) be its
completion. All finite stages are actual embedded fields
\(K\subseteq E\subseteq M\), complete with their prescribed absolute
values. Write \(p_E:X_E\to X\), \(q_E:X_L\to X_E\), and
\(\pi:X_L\to X\). Scalar extensions use actual completed affinoid algebras.

\textbf{Theorem A (exact compact tests).} If \(X\) is good, separated
and strict, then for every pair of nonnegative integers \((a,b)\) and
every ordinary CLD form \(\beta\in A_c^{a,b}(X_L)\), there are one
finite \(E\subseteq M\) and \(\alpha_E\in A_c^{a,b}(X_E)\) such that
\(\beta=q_E^*\alpha_E\) as smooth sections. Consequently both
differentials, their compositions and exterior products with pulled-back
sections respect this identity. No dimension, boundarylessness or
positivity is required. This is descent of a compact smooth form, rather
than of an arbitrary analytic function or an arbitrary upstairs space
presentation.

\textbf{Theorem B (the new gate).} Suppose in addition that X is
boundaryless and pure of analytic dimension n\ensuremath{\ge}1. Retain its possibly
singular and nonreduced analytic structure. Then X\_L has these same
ambient properties and is pure of dimension n.~For every finite
continuous h:X\ensuremath{\to}R,

\[
 d'd''[h]=0\quad\Longleftrightarrow\quad d'd''[h\circ\pi]=0.       \tag{T}
\]

Both equations are tested against \textbf{all} compact ordinary smooth
(n\ensuremath{-}1,n\ensuremath{-}1) forms and use the canonical CLD calibration for their own
fields. All assertions remain true on every ordinary open subspace. No
separability, tameness, perfectness, unit-separation, Radon-Hessian,
psh, or invariant-form-descent assumption is imposed. Nilpotents and all
actual local fibre lengths are retained.

The theorem does not recognize h as classically pluriharmonic. It does
not prove two-sign regularization, transport along arbitrary curve maps,
or transport through a general transcendental analytic extension. In
particular L/K may fail to be algebraic, but here it is the completion
of the specified algebraic extension. The original P06.2 remains open;
P05.2 is unchanged. The full input is preserved in original-task.json,
original-problem.txt, completion-criteria.txt, and
assignment-verbatim.txt.

All opens and supports use the ordinary Hausdorff topology. An affinoid
neighborhood contains an ordinary neighborhood. Its ordinary interior
inside X is denoted int\_X(V); this differs from its absolute analytic
interior \(\operatorname{Int}(V)\). Set \(D=d'd''\) and \(t_f=-\log|f|\) for an analytic unit. Up to the fixed current-duality sign, which has no effect
on vanishing, the equation is

\[
 B_h(\alpha):=\int_X h\,(D\alpha)_{\mathrm{Rad},K}=0
 \quad(\alpha\in A_c^{n-1,n-1}(X)).                         \tag{1}
\]

Only smooth differentiated top forms are assigned Radon measures here.
The current D{[}h{]} is not assumed to be a measure.

\subsection{2. Exact primary foundations and their
applicability}\label{exact-primary-foundations-and-their-applicability}

CLD means Chambert-Loir--Ducros, \emph{Formes différentielles réelles et
courants sur les espaces de Berkovich}, arXiv:1204.6277v2, 25 July 2025.
The complete selected primary statements and their printed proof bodies
are embedded below. These are imported analytic foundations, not
previous campaign claims. The algebra, support, collar, and finite-stage
arguments of Sections 3--10 are supplied in this document.

\textbf{Affinoid foundations.} CLD 11.1.1--2: affinoid spectra are
compact Hausdorff; affinoid algebras are noetherian with closed ideals;
their analytic spectra surject onto their prime spectra. Completed
residue fields are completions of the fraction fields of the
point-kernel quotients. In particular an element with no analytic zero
is an actual unit, including in a nonreduced affinoid algebra. CLD
11.1.5: a rational system \textbar f\_i\textbar\ensuremath{\le}r\_i\textbar g\textbar{}
with positive real r\_i and no common analytic zero defines an affinoid
domain, and every point has such a neighborhood inside any prescribed
ordinary open, with all inequalities strict at the point. Its algebra is
the quotient obtained by adjoining the bounded ratios f\_i/g. CLD
11.1.7--10: good spaces have affinoid neighborhood bases; opens are
domains; domains have the universal analytic factorization property;
scalar affinoids use completed tensors and scalar projection is
surjective. CLD 11.4.4 gives finite algebras their finite-module
affinoid topology. CLD 11.7.7 gives strict affinoid neighborhoods in
good strict spaces over a nontrivially valued field. Applying this to
ordinary opens gives a strict neighborhood inside any such open. We use
these facts on actual algebras, never spectral norms on a reduction as
substitute Banach norms.

\textbf{Dimensions and boundaries.} CLD 11.2.2--3: separatedness implies
Hausdorffness; strict affinoid dimension is Krull dimension; complete
scalar extension preserves total analytic dimension; domains preserve
local dimension. For the minimal components V\_i of an affinoid V,
dim\_y(V)=max\_\{y\ensuremath{\in}V\_i\}dim(V\_i). Purity means constant local
dimension. These statements alone do not assert purity after completion.
Sections 3 and 9 prove that separate assertion. CLD 11.4.1 gives the
relative interior composition equality when the middle morphism is
locally separated. Berkovich, \emph{Étale cohomology for non-Archimedean
analytic spaces}, Publ. Math. IHES 78 (1993), Proposition
1.5.5(i),(iii),(iv), gives the ordinary relative interior of a domain
and the base-change inclusion for relative interiors, including complete
field extension. Thus ordinary opens in a boundaryless space and scalar
extensions of a boundaryless space are boundaryless. Equality of
arbitrary absolute interiors under scalar extension is not used.

\textbf{Ordinary forms and nullity.} CLD 6.2.1--2 defines ordinary
smooth forms by tropical descriptions on ordinary open neighborhoods and
gives ordinary sheaves closed under wedge, J,d',d'\,'. Proposition
15.1.1(b) gives descriptions by finitely many analytic units on an
ordinary open cover, in every bidegree. CLD 3.3.3 and Lemma 1.9.4 give
ambient Euclidean coefficient representatives, using the latter only
when the paralinear subset is closed in the ambient open. CLD 3.3.10--12
gives the exterior calculus and functorial smooth pullback. Proposition
3.4.4 states exactly: for a tropical chart f and coefficient form \ensuremath{\omega},
f\^{}*\ensuremath{\omega} is zero in the G-form sheaf iff \ensuremath{\omega} restricts to zero on f(V) for
every compact domain V. We use this to identify scalar sections that are
equal as functions and to identify forms in common moment charts.
Ordinary forms are sections of that sheaf with ordinary local
descriptions; no descent of arbitrary invariant forms is imported.

\textbf{Compact cutoffs.} CLD 15.2.2 makes good Hausdorff spaces locally
holomorphically separated. Corollary 15.2.5 says that, for a compact C
inside an ordinary open U in such a space, a smooth function exists
which is one near C and has compact support contained in U. Section 3
supplies the finite compact partition and an independent cutoff
construction. The overly broad fixed-index compact-support reading of
15.2.9 is not used.

\textbf{Polyhedra and measures.} Ducros, \emph{Espaces de Berkovich,
polytopes, squelettes et théorie des modèles}, Confluentes Math. 4
(2012), 1250007, Theorem 3.2(1), states: if a compact \ensuremath{\Gamma}-strict
k-analytic space has dimension d and f is a finite tuple of analytic
units, \textbar f\textbar(X) is a c-polytope of dimension at most d,
with c=(Q,(\textbar k\^{}*\textbar \ensuremath{\Gamma})\^{}Q). We apply this to strict
compact affinoids, \ensuremath{\Gamma}=\{1\}; it permits nilpotents and singularities. No
universal dimension equality from CLD 12.2.4 is used. CLD 5.4.4 and
14.1.8 make compact charts pleasant; 14.1.12 gives compatible skeletal
and image cell refinements. CLD 14.2.2 and 14.3.2--5 give positive
constant canonical calibrations on n-cells: the weight is the actual
finite-flat moment degree times a primitive volume vector. CLD 14.5.1--2
sums these weights on image cells. CLD 5.8.1--5 constructs top-form
Radon measures from their skeletal densities, compatibly on domains;
5.10.1(c),(d) gives the corresponding pleasant-chart integration
formula. Smooth top forms are locally integrable and compact ones
finite. Lower-dimensional image cells have no n-dimensional top measure.
Section 5 derives the variation bound from the positive sheet densities,
rather than from a signed pushforward equality.

\textbf{Two different scalar comparisons.} CLD Proposition 14.3.6 states
\ensuremath{\pi}\_*\ensuremath{\mu}\_L=\ensuremath{\mu} for canonical skeletal calibrations under a complete
extension. Proposition 15.3.1 states that a locally integrable top
G-form \ensuremath{\omega} satisfies

\[
 \pi_*((\pi^*\omega)_{\mathrm{Rad},L})
                 =\omega_{\mathrm{Rad},K}.                 \tag{2}
\]

The Borel variation measures have the analogous identity. There is no
field-degree factor. CLD Remark 14.4.3 instead states: on the
\textbf{same carrier}, an F-analytic space viewed as K-analytic, for
finite F/K, has the same dimension and skeletal subsets, and its
calibration is

\[
                      \mu_K=[F:K]\mu_F.                    \tag{3}
\]

We use (3) only after constructing the actual local F-analytic structure
and identifying the forms. Equation (2) is used only on smooth top
forms, followed by a continuous multiplier, not as continuous-current
pullback. Its primary proof uses CLD 13.3.3, the degree-sum formula for
actual Artinian fibre algebras; Section 10 gives that algebra
explicitly.

\textbf{Currents and Stokes.} CLD 7.3.5 defines {[}h{]} for a locally
integrable multiplier, including finite continuous functions. CLD
Theorem 15.7.1 gives Hausdorff pure n-dimensional analytic spaces their
level-n tropical structure and analytic boundary; 15.7.2 gives domain
compatibility. Stokes, CLD 5.11.1(a), applies to compact C\^{}1 (n\ensuremath{-}1,n)
forms; J gives the double-prime version. On our boundaryless spaces both
boundary integrals vanish. Section 5 uses these twice on arbitrary
complementary forms; the symmetry-restricted Green formula is not used.

\subsection{3. Actual finite algebras, geometry, and compact
partitions}\label{actual-finite-algebras-geometry-and-compact-partitions}

\subsubsection{3.1. Finite topology and elementary field
facts}\label{finite-topology-and-elementary-field-facts}

A finite-dimensional normed vector space over a complete valued field
has a norm equivalent to the coefficient maximum norm in every basis.
Induct on dimension: the preceding span is complete and closed, so the
last basis vector has positive distance from it. Dividing a vector by
its last coefficient bounds that coefficient by its norm; the triangle
inequality and induction bound the remaining coefficients. The reverse
bound is the triangle inequality. Consequently, for any affinoid Banach
algebra A and finite E/K of degree m, tensoring the bounded coordinate
maps E\ensuremath{\to}K identifies the tensor topology with that of A\^{}m. The latter
is complete. Hence

\[
            A\widehat\otimes_K E=A\otimes_K E               \tag{4}
\]

with its actual finite-module topology, including nilpotents.

A finite field extension is complete for its extending norm; existence
follows by giving the finite algebra its affinoid finite-module topology
and taking a point of its nonempty spectrum. Its kernel is zero. Any two
extending absolute values are equivalent finite-dimensional norms; apply
their comparison constants to powers z\^{}j and take j-th roots to prove
equality. Thus all finite field embeddings here preserve the norm.

The elementary Galois facts needed below can be reconstructed as
follows. A finite separable extension has degree-many embeddings into an
algebraic closure, by successively choosing the distinct roots along a
finite generating tower. K is infinite, since a finite field has no
nontrivial absolute value. For generators a\_1,\ldots,a\_s choose
coefficients (1,t,\ldots,t\^{}\{s\ensuremath{-}1\}) such that their linear
combination has distinct values under all embeddings: each pair forbids
the roots of one nonzero polynomial in t, and finitely many forbidden
roots cannot exhaust K. The combination is primitive by comparing the
number of its conjugates with the degree. In a finite normal separable
extension E/K all embeddings of intermediate fields extend to E by
adjoining successive roots and normality; the image is again E. Thus
E\^{}G=K. For H\ensuremath{\le}G and F=E\^{}H, choose a K-primitive element b for E.
Then E=F{[}b{]}, and its H-orbit polynomial has coefficients in F. The
F-minimal polynomial divides this polynomial and has all its
\textbar H\textbar{} distinct roots, so its degree is
\textbar H\textbar. Thus E/F is Galois with group H. Conjugating
embeddings of F conjugates H. A splitting field of a separable primitive
polynomial is the actual finite Galois normal closure used in Section 6.

For a finite Galois E/K, tensoring the finite kernel equations defining
E\^{}G=K with the K-vector space A gives

\[
 (A\otimes_K E)^G=A,\qquad
 N(f):=\prod_{g\in G}gf\in A^\times\quad(f\in(A\otimes E)^\times). \tag{5}
\]

The inverse product descends as well. No averaging in K is used. The
automorphisms are bounded on the finite-module algebra. Its restricted
K-spectrum and E-spectrum agree because all bounded multiplicative
seminorms restrict to the unique norm on E. Their actual rational
algebras and unit germs agree; the action on points satisfies
\textbar gf(y)\textbar=\textbar f(g\^{}\{-1\}y)\textbar.

\subsubsection{3.2. All minimal components at finite
stages}\label{all-minimal-components-at-finite-stages}

In a strict affinoid domain M(A) in pure n-dimensional X, every minimal
prime P has \(\dim(A/P)=n\). Spectral surjectivity supplies a point of
kernel P; it lies on that component alone, and the primary local
dimension formula gives this equality.

First let E/K be finite separable. For the nilradical J of A, J is
nilpotent, and the injection A/J\ensuremath{\to}\ensuremath{\prod}\_P Frac(A/P) stays injective under
algebraic tensor over K. Each field tensor with E is a finite product of
fields (factor a separable primitive polynomial). Thus (A/J)\ensuremath{\otimes}E is
reduced, and the nilradical of A\ensuremath{\otimes}E is precisely J\ensuremath{\otimes}E. The finitely many
maps of this reduced algebra into those field factors have prime kernels
with zero intersection. A minimal prime Q contains their product, hence
one kernel, and by minimality equals it. Its contraction is the
associated P. The finite integral injection A/P\ensuremath{\to}(A\ensuremath{\otimes}E)/Q gives equal
Krull dimensions.

Here are the elementary integral dimension facts used in that last step.
For a finite integral injection of domains C\ensuremath{\to}D, localization at a prime
of C is a nonzero finite module. Its reduction modulo the maximal ideal
is nonzero by the determinant argument for Nakayama: if generators equal
a maximal-ideal-valued matrix times themselves, det(I\ensuremath{-}matrix), a unit,
annihilates them. A maximal ideal of the finite quotient lies over the
given prime (a subring over which a field is integral is a field).
Applying lying over after a quotient gives going up. For
incomparability, quotient by the smaller upstairs prime and invert the
nonzero elements of the base domain. The resulting finite integral
domain over a field is a field, since a least-degree monic equation
expresses the inverse of any nonzero element. A nonzero element of the
quotient domain cannot disappear in that localization. Thus strictly
nested primes cannot have the same contraction. Chain contraction and
chain lifting prove equal dimensions. This argument needs no
completeness of the fraction fields.

For finite purely inseparable E/K in characteristic p, choose q=p\^{}r
with E\^{}q\ensuremath{\subseteq}K. In the actual algebra B=A\ensuremath{\otimes}E every b satisfies b\^{}q\ensuremath{\in}A.
A\ensuremath{\to}B is an injection, finite and integral. For each prime P of A there is
exactly one prime Q above it: membership b\ensuremath{\in}Q is equivalent to b\^{}q\ensuremath{\in}P.
Lying over gives existence. This formula also preserves and reflects
prime inclusions, so minimal primes correspond and dimensions of their
quotients agree by the preceding integral argument. Every component of B
has dimension n, whether or not B is reduced. Combining the two cases in
the embedded mixed tower of Section 6 proves finite-stage purity for
every finite algebraic E/K.

\subsubsection{3.3. Other ambient properties and compact scalar inverse
images}\label{other-ambient-properties-and-compact-scalar-inverse-images}

For any complete extension, scalar affinoids are actual completed
tensors. Inverse images of downstairs affinoid neighborhoods contain
inverse images of their ordinary interiors, so goodness is preserved.
Strict presentations base change to strict presentations. Closed
diagonal base change proves separatedness; hence Hausdorffness.
Relative-interior base-change inclusion proves boundarylessness.
Ordinary opens retain all these properties. Finite-stage purity has just
been proved; completed-stage purity is established independently in
Section 9.

Scalar projection has compact inverse images. Cover a downstairs compact
C by finitely many ordinary interiors of affinoid neighborhoods V\_i.
Its inverse image is closed in the finite union of the compact
(V\_i)\_L, because the carrier is Hausdorff. This applies also to an
infinite completed extension. No finite-fibre assertion is used for that
extension.

We also need a finite full-fibre shrinking argument. Inside an affinoid
neighborhood V\_0 of x, if an ordinary open N\ensuremath{\subseteq}(V\_0)\_E contains the
whole fibre, its compact complement has closed image excluding x. A
smaller strict affinoid neighborhood V of x can be chosen in the
complement of that image and in int\_X(V\_0). Then V\_E\ensuremath{\subseteq}N.

\subsubsection{3.4. Valid cutoffs, finite partitions, and continuous
approximation}\label{valid-cutoffs-finite-partitions-and-continuous-approximation}

The compact-cutoff input can also be proved directly. Given x\ensuremath{\in}U ordinary
open, choose an affinoid neighborhood V of x and let
C\_0=V(int\_X(V)\ensuremath{\cap}U). It is compact and excludes x. For each y\ensuremath{\in}C\_0 an
analytic function on V separates x,y by its norm. A smooth real function
of its logarithm can be chosen one near the first value and zero near
the second. If either norm is zero, make the composing function constant
on a sufficiently negative half-line and extend by that constant across
the zero locus. This extension is locally constant there and ordinary
smooth; elsewhere the analytic function is locally a unit. A finite
product of these functions is one near x and zero on a relative
neighborhood of C\_0. Its zero extension outside V is smooth: on
int\_X(V) and outside V this is immediate, and at C\_0 relative
vanishing gives an ordinary zero neighborhood. Its support is compact
inside U. Finite sums of squares and smooth composition give the same
assertion for compact C\ensuremath{\subseteq}U.

For a finite open cover U\_i of compact C, choose finitely many such
nonnegative bumps with positive loci covering C, grouped as \ensuremath{\psi}\_i
supported compactly inside U\_i. Put s=\ensuremath{\sum}\ensuremath{\psi}\_i, and choose compact \ensuremath{\chi}=1
near C supported inside \{s\textgreater0\}. The functions \ensuremath{\theta}\_i=\ensuremath{\chi}\ensuremath{\psi}\_i/s
on \{s\textgreater0\}, zero elsewhere, are ordinary smooth with compact
support inside U\_i and sum one near C. Zero extension is smooth because
each support is a compact closed subset of its defining open. This
construction requires no global paracompactness. Applying it to
small-oscillation neighborhoods of a continuous h gives a compact
ordinary smooth a with \textbar a\ensuremath{-}h\textbar\textless \ensuremath{\varepsilon} on any fixed
compact C.

\subsection{4. Reconstructed finite-Galois supporting
theorem}\label{reconstructed-finite-galois-supporting-theorem}

This section proves the finite-Galois case of (T), independently of the
round-19 conclusion. The stabilizer, Newton, tensor, support, and collar
steps are included because a scalar-test-only finite theorem would not
support the completion argument.

\subsubsection{4.1. Actual fibres and
stabilizers}\label{actual-fibres-and-stabilizers}

For x\ensuremath{\in}X put C=H(x). The actual fibre is M(C\ensuremath{\otimes}\_K E). Factoring a
separable primitive polynomial gives a finite product of fields D\_y,
each with a single analytic point by finite norm uniqueness. G acts
transitively on those factors: a union of two distinct orbits would give
a nontrivial invariant idempotent, whereas (C\ensuremath{\otimes}E)\^{}G=C by the kernel
calculation. For the actual point stabilizer H, invariant tuples
identify D\_y\^{}H=C. The image of the primitive b has
\textbar H\textbar{} distinct conjugates in D\_y and D\_y=C{[}b{]}. The
same orbit-polynomial argument gives {[}D\_y:C{]}=\textbar H\textbar{}
and group H. In particular E\^{}H embeds in C. If y is not G-fixed, H is
proper. This uses actual tensor factors, not a covering-space model.

\subsubsection{4.2. A separable residue constant lifts by
quotient-Banach
Newton}\label{a-separable-residue-constant-lifts-by-quotient-banach-newton}

Suppose a specified finite separable F/K embeds in H(x). Choose a
primitive a with monic minimal polynomial P(T)=\ensuremath{\sum}p\_jT\^{}j. In a strict
affinoid neighborhood of x, fractions of the kernel quotient are dense
in H(x). Approximate a by g=f/b, retaining b invertible on a strict
rational neighborhood where \textbar b\textbar\ensuremath{\ge}\textbar k\_0\textbar{}
is strictly satisfied at x. Choose R\ensuremath{\in}\textbar K\^{}\emph{\textbar{} with
R\textgreater max(1,\textbar a\textbar), d\ensuremath{\in}\textbar K\^{}}\textbar{}
with 0\textless d\textless\textbar P'(a)\textbar, and

\[
 M_0=\max\bigl(\{1\}\cup\{|p_j|R^{j-2}:2\le j\le\deg P\}\bigr).
\]

Choose c\ensuremath{\in}\textbar K\^{}*\textbar{} so that s=c/d\textless min(1,R) and
v=M\_0c/d\^{}2\textless1. Approximate closely enough that
\textbar g(x)\textbar\textless R, \textbar P(g)(x)\textbar\textless c,
\textbar P'(g)(x)\textbar\textgreater d, and
\textbar g(x)\ensuremath{-}a\textbar\textless s. Retain the strict rational
neighborhood \textbar g\textbar\ensuremath{\le}R, \textbar P(g)\textbar\ensuremath{\le}c,
\textbar P'(g)\textbar\ensuremath{\ge}d.~Its actual quotient Tate algebra has
\textbar\textbar g\textbar\textbar\ensuremath{\le}R,
\textbar\textbar P(g)\textbar\textbar\ensuremath{\le}c,
\textbar\textbar P'(g)\textsuperscript{\{-1\}\textbar\textbar\ensuremath{\le}d}\{-1\},
by the adjoined quotient variables. These are Banach bounds imposed by
relations, valid with nilpotents; spectral smallness is not being used.

Binomial expansion, with integer coefficients of norm at most one, gives
for \textbar\textbar z\textbar\textbar\ensuremath{\le}R and
\textbar\textbar w\textbar\textbar\ensuremath{\le}min(1,R)

\[
 \|P(z+w)-P(z)-P'(z)w\|\le M_0\|w\|^2,
 \qquad\|P'(z+w)-P'(z)\|\le M_0\|w\|.
\]

Iterate
\[
 z_{j+1}=z_j-P'(z_j)^{-1}P(z_j),\qquad z_0=g.
\]
The first increment has norm at most \(s\). Neumann inversion retains
its derivative inverse bound because \(M_0s/d=v<1\). The next
increment is at most
\[
 (M_0/d)\|z_{j+1}-z_j\|^2\le v\|z_{j+1}-z_j\|.
\]
All \(z_j\) have norm at most \(R\), so they converge in the actual
algebra to \(z\), with \(P(z)=0\) and \(\|z-g\|\le s\).
Two roots within \(s\) of \(g(x)\) have difference multiplied by
\(P'(g(x))+w\), where \(|w|\le M_0s<d\), equal to zero.
This nonzero factor proves uniqueness, hence \(z(x)=a\).
The map \(F=K[T]/P\to\mathcal{O}(V)\), \(T\mapsto z\), is
injective (a unital map from a field to a nonzero algebra) and bounded
by finite norm equivalence.

This gives an actual \(F\)-analytic neighborhood. For the finite
\(F\)-action on \(A=\mathcal{O}(V)\), \(\|u\|_A\le C|u|\).
The equivalent complete norm
\[
 \|b\|_F=\sup_{u\in F^*}\frac{\|ub\|_A}{|u|}
\]
is \(F\)-homogeneous and submultiplicative:
\[
 \|ab\|_F\le\|a\|_F\|b\|_A\le\|a\|_F\|b\|_F.
\]
A strict K-Tate presentation has uniformly bounded monomials also in
this equivalent norm and thus extends to a bounded surjection F\{T\}\ensuremath{\to}A.
The quotient norm is equivalent: boundedness gives one bound, and
uniform lifts through the original K-Tate presentation give the other.
The spectra coincide by finite norm uniqueness, and actual rational
algebras coincide by their universal presentations. These rational bases
identify analytic sheaves on the common ordinary carrier. Common unit
moments and compact-image nullity identify all ordinary forms,
operations, and supports over F and K.

Restrict to an ordinary U around x. It is good, strict, separated and
pure over F: its affinoid algebras and Krull component dimensions are
unchanged. The F-diagonal is the closed equalizer inside the K-fibre
product, so separatedness is retained. Boundarylessness over F follows
from interior composition for U\ensuremath{\to}M(F)\ensuremath{\to}M(K), with finite separated second
map. On this actual same carrier, (3) gives for every compact
complementary test

\[
                         B_{h,K}=[F:K]B_{h,F}.             \tag{6}
\]

The positive real integer is cancellable even if its image in K is zero.

\subsubsection{4.3. Strong induction off the fixed
locus}\label{strong-induction-off-the-fixed-locus}

Induct simultaneously for every field and space on \textbar G\textbar.
The identity case is immediate. Let S be the G-fixed locus in X\_E,
closed by Hausdorffness and the finite equalizer description. At y\ensuremath{\notin}S,
the proper stabilizer H gives a specified F=E\^{}H\ensuremath{\subseteq}H(x). Section 4.2
lifts it on U; (6) shows zero F-current there. The actual tensor
identities are

\[
 F\otimes_K E\simeq\prod_{\tau:F\hookrightarrow E}E,
 \quad A\otimes_K E\simeq
             \prod_{\tau}A\otimes_{F,\tau}E.                \tag{7}
\]

They follow from the distinct linear factors of a primitive polynomial
and the Chinese remainder theorem, and retain all nilpotents of A.
Finite-module topologies identify their analytic components, which are
open and closed. Each is scalar extension by E/\ensuremath{\tau}(F), a Galois extension
with group conjugate to H, of smaller order. All induction hypotheses
were checked above. Thus the current vanishes near y. Compact partitions
show that it vanishes on every complementary test supported outside S.

\subsubsection{4.4. Norm units match arbitrary tensor tests on the fixed
subspace}\label{norm-units-match-arbitrary-tensor-tests-on-the-fixed-subspace}

Let \ensuremath{\beta} be any compact ordinary complementary test. At y\ensuremath{\in}S\ensuremath{\cap}supp(\ensuremath{\beta}),
transitivity makes its entire fibre the singleton \{y\}. A local
description uses units f\_1,\ldots,f\_k and a Euclidean smooth superform
\ensuremath{\omega}. Intersect the G-translates of the description/unit open. Full-fibre
shrinking produces a strict affinoid V downstairs with V\_E inside it.
Include all gf\_i in the orbit moment t. G permutes its coordinates. Let
W be the fixed linear subspace, i its inclusion, and
P\_0=\textbar G\textbar\^{}\{-1\}\ensuremath{\sum}\_g g the real averaging projection
onto W. At actual fixed points t\ensuremath{\in}W, and its averaged orbit coordinates
are

\[
       |G|^{-1}\sum_g t_{gf_i}
          =|G|^{-1}t_{N(f_i)}\circ p_E.                     \tag{8}
\]

All norms in (8) are actual downstairs units by (5); division is real.
Shrink V so that both t(V\_E) and iP\_0t(V\_E) lie in the Euclidean
coefficient open. This is possible because both agree at the singleton
fibre and full-fibre shrinking preserves these open conditions. Extend \ensuremath{\omega}
globally by a Euclidean bump unchanged near those compact images. Set
\ensuremath{\omega}\_0=P\_0\^{}\emph{i\^{}}\ensuremath{\omega}. Then i\textsuperscript{\emph{(\ensuremath{\omega}\ensuremath{-}\ensuremath{\omega}\_0)=0 and
t\^{}}\ensuremath{\omega}\_0=p\_E}*\ensuremath{\alpha}\_0 for a downstairs ordinary complementary form
obtained using (8). This matches on W; it does not descend \ensuremath{\beta} or an
arbitrary invariant form.

Cover p\_E(S\ensuremath{\cap}supp \ensuremath{\beta}) by finitely many such ordinary downstairs opens,
and take compact \ensuremath{\theta}\_i supported there with sum one near that compact.
Decompose \ensuremath{\beta} into \ensuremath{\beta}\_rem=(1\ensuremath{-}p\_E\^{}\emph{\ensuremath{\sum}\ensuremath{\theta}\_i)\ensuremath{\beta} and \ensuremath{\beta}\_i=p\_E\^{}}\ensuremath{\theta}\_i
\ensuremath{\beta}. The remainder is supported outside S. Put
\ensuremath{\beta}\_\{0,i\}=p\_E\^{}*(\ensuremath{\theta}\_i\ensuremath{\alpha}\_\{0,i\}), \ensuremath{\gamma}\_i=\ensuremath{\beta}\_i\ensuremath{-}\ensuremath{\beta}\_\{0,i\}, with all
products extended by zero. Their supports lie in fixed compact inverse
images of supp \ensuremath{\theta}\_i inside the charts. Equation (2) gives
B\_\{h\_E\}(\ensuremath{\beta}\_\{0,i\})=B\_h(\ensuremath{\theta}\_i\ensuremath{\alpha}\_\{0,i\})=0. Each \ensuremath{\gamma}\_i has
coefficient restriction zero on W. Coordinates used by \ensuremath{\theta}\_i enlarge this
to W times their entire coordinate space. This property persists under
any further moment refinement.

\subsection{5. Tensor collars, actual variation, and continuous
multipliers}\label{tensor-collars-actual-variation-and-continuous-multipliers}

\textbf{Collar lemma.} Let Q be a fixed linear map with kernel W. Let \ensuremath{\sigma}
be a compact n-dimensional affine polytope and \ensuremath{\xi} a smooth (n\ensuremath{-}1,n\ensuremath{-}1)
superform defined near \ensuremath{\sigma} and its nearby W comparison points, with
i\^{}*\ensuremath{\xi}=0 on W. Let \ensuremath{\chi} be smooth radial, one near zero and zero outside
the unit ball, and set \ensuremath{\xi}\_\ensuremath{\delta}=\ensuremath{\chi}(Qz/\ensuremath{\delta})\ensuremath{\xi}. On the tangent of \ensuremath{\sigma}, the
coefficient L\^{}1 norm of D\ensuremath{\xi}\_\ensuremath{\delta} is bounded uniformly for small \ensuremath{\delta}; the
L\^{}1 norm of \ensuremath{\xi}\_\ensuremath{\delta} tends to zero. Extra coordinates annihilated by Q
are allowed.

\textbf{Proof.} Let H be the tangent and r=rank(Q\textbar H). If the
affine span misses W, its Q-image is a closed affine space at positive
distance from zero, so the collar is eventually empty. Otherwise choose
w\_0 in that span\ensuremath{\cap}W and write z=w\_0+u+v, with u\ensuremath{\in}H\ensuremath{\cap}W and v in a fixed
r-dimensional complement. The map Q on the complement is injective;
\textbar Qz\textbar\ensuremath{\le}\ensuremath{\delta} bounds v in an O(\ensuremath{\delta}) box, while compactness bounds
u in a fixed box. Fubini with the fixed Jacobian bounds collar volume by
O(\ensuremath{\delta}\^{}r). The comparison point w\_0+u need not be in \ensuremath{\sigma}, which handles
contact only on a boundary face. If r=0 the cell either misses the
collar or lies in W, on which \ensuremath{\xi} and its restricted derivatives vanish.
If r\ensuremath{\ge}2, cutoff derivatives through order two are O(\ensuremath{\delta}\^{}\{-2\}); bounded
derivatives of \ensuremath{\xi} times O(\ensuremath{\delta}\^{}r) volume give the asserted bound.

If r=1 use tangent coordinates (u\_1,\ldots,u\_\{n\ensuremath{-}1\},t). The
restricted cutoff depends only on t. In its second-derivative term
\ensuremath{\chi}'`(t)d't\ensuremath{\wedge}d''t\ensuremath{\wedge}\ensuremath{\xi}, every exterior component with either transverse
generator is killed by repetition. The unique surviving basis component
is

\[
 A(u,t)\,d'u_1\wedge\cdots\wedge d'u_{n-1}
              \wedge d''u_1\wedge\cdots\wedge d''u_{n-1}.
\]

Restriction-zero gives A(u,0)=0, so smoothness gives
\textbar A(u,t)\textbar\ensuremath{\le}C\textbar t\textbar{} on the fixed compact
comparison region. The second-cutoff term is therefore O(\ensuremath{\delta}\^{}\{-1\});
first-cutoff terms also are O(\ensuremath{\delta}\^{}\{-1\}). O(\ensuremath{\delta}) collar volume gives
uniform integrability. This includes n=1. Finally \ensuremath{\xi} is bounded, so its
undifferentiated L\^{}1 norm tends to zero on positive-rank collars;
rank-zero restrictions vanish. Finite cells and fixed positive weights
preserve both conclusions. This proves the lemma.

The lemma controls the variation of differentiated \textbf{smooth} top
forms on actual carriers. On a strict compact affinoid with fixed moment
t, the cited polyhedral and skeletal refinements give finitely many
n-sheets above each image n-cell. Each has constant positive weight, its
actual finite-flat fibre degree times a primitive volume. Thus for any
represented top coefficient form \ensuremath{\zeta},

\[
 \|(t^*\zeta)_{\mathrm{Rad}}\|_{\mathrm{TV}}
 \le\sum_{\sigma}c_\sigma
        \int_\sigma\|\zeta|_{T\sigma}\|\,d\mathrm{vol}_\sigma,  \tag{9}
\]

with fixed finite constants. Indeed evaluate the scalar density
separately on every sheet, take its absolute value there, and sum the
positive weights. This is the density construction in 5.8.4--5 and
14.3.2; it cannot be replaced by cancellation in a signed pushforward.
All nilpotent lengths enter the positive weights. Lower-dimensional
moment images give zero top form. Covering a fixed compact support by
finitely many strict compact charts and overcounting their positive
variations gives a valid bound; no partition derivatives are needed for
this estimate.

Fix one \ensuremath{\gamma}=\ensuremath{\gamma}\_i from Section 4.4, with Q=1\ensuremath{-}P\_0. Define e\_\ensuremath{\delta}=\ensuremath{\chi}(Qt/\ensuremath{\delta})\ensuremath{\gamma},
extended by zero in its chart. Its support lies in one fixed compact C.
The test (1\ensuremath{-}\ensuremath{\chi}(Qt/\ensuremath{\delta}))\ensuremath{\gamma} is supported outside S: Qt=0 on S and \ensuremath{\chi} is one
near zero. The induction therefore gives the exact identity

\[
                  B_{h_E}(\gamma)=B_{h_E}(e_\delta).        \tag{10}
\]

On a finite strict compact cover of C, add the moment coordinates
describing the fixed \ensuremath{\theta}\_i. The coefficient of \ensuremath{\gamma} still restricts to zero
on W times the added space. On each chart extend the smooth coefficient
of \ensuremath{\theta}\_i to the full ambient Euclidean space by a bump unchanged near its
compact image; multiplying the globally defined \ensuremath{\omega}\ensuremath{-}\ensuremath{\omega}\_0 by this extension
preserves restriction-zero on the whole enlarged W. This ensures that
the off-cell comparison points in the collar proof lie in a valid smooth
coefficient domain. The collar lemma and (9) give

\[
     \sup_{0<\delta<\delta_0}
       \|(De_\delta)_{\mathrm{Rad},E}\|_{\mathrm{TV}}=C_1<\infty. \tag{11}
\]

For any fixed ordinary smooth compact scalar a, two Stokes applications
on the pure boundaryless space give

\[
                   \int a\,De_\delta
                       =\int e_\delta\wedge Da.             \tag{12}
\]

Explicitly, Stokes on a d'`e gives \ensuremath{\int}aDe=\ensuremath{-}\ensuremath{\int}d'a\ensuremath{\wedge}d'`e. Stokes on d'a\ensuremath{\wedge}e
gives \ensuremath{\int}d'a\ensuremath{\wedge}d''e=\ensuremath{-}\ensuremath{\int}Da\ensuremath{\wedge}e; e has even total degree, so Da\ensuremath{\wedge}e=e\ensuremath{\wedge}Da. No
symmetry of the complementary form is required. To estimate the right
side, add \textbf{a's moment coordinates} on a finite compact cover. The
enlarged fixed space is W times the entire extra space. Restriction-zero
persists; extend the local scalar coefficients by Euclidean bumps near
each compact joint image as above. The undifferentiated collar lemma on
the joint moment gives L\^{}1 convergence of e\_\ensuremath{\delta}, and Da has bounded
coefficients on these compact images. Formula (9) for e\_\ensuremath{\delta}\ensuremath{\wedge}Da proves
\ensuremath{\int}a(De\_\ensuremath{\delta})\_Rad\ensuremath{\to}0. If the original moment has dimension \textless n,
De\_\ensuremath{\delta} is zero; that observation alone is not substituted for this
joint-moment estimate. The weak estimate may depend on the fixed a,
whereas (11) was established independently of a.

Choose a with \textbar a\ensuremath{-}h\_E\textbar\textless \ensuremath{\varepsilon} on C using Section 3.4.
Equation (11) yields \[
 \limsup_{\delta\to0}\left|\int h_E(De_\delta)_{\mathrm{Rad}}\right|
                                      \le C_1\varepsilon.
\] Let \ensuremath{\varepsilon}\ensuremath{\to}0, and use (10). Thus B\_\{h\_E\}(\ensuremath{\gamma})=0. The norm-matched tests,
\ensuremath{\gamma}\_i, and off-S remainder exhaust the original arbitrary \ensuremath{\beta}, proving
finite-Galois ascent. Reverse transport follows at once from (2) applied
to D\ensuremath{\alpha} for each downstairs compact \ensuremath{\alpha}, whose pullback is compact by
Section 3.3. This finishes the reconstructed finite-Galois supporting
theorem.

\subsection{6. Actual mixed finite stages and purely inseparable higher
forms}\label{actual-mixed-finite-stages-and-purely-inseparable-higher-forms}

\subsubsection{6.1. Finite separable transport uses an actual normal
closure}\label{finite-separable-transport-uses-an-actual-normal-closure}

For finite separable E/K, choose its actual finite Galois normal closure
N/K in an algebraic closure containing E, with the unique extending
norm. Finite purity and all ambient properties hold. Finite-Galois
ascent makes h\_N zero-current. For each compact complementary test \ensuremath{\alpha}\_E
on X\_E, its compact pullback to X\_N has pairing equal to
B\_\{h\_E\}(\ensuremath{\alpha}\_E) by (2). Thus h\_E has zero current. This is descent
through N/E by the canonical top measure, not a normalization by
{[}N:E{]} or an abstract degree argument. The construction of N does not
require N\ensuremath{\subseteq}M. Reverse finite transport uses (2) directly.

\subsubsection{6.2. Finite purely inseparable all-bidegree
descent}\label{finite-purely-inseparable-all-bidegree-descent}

For finite purely inseparable E/K of exponent q=p\^{}r and actual
affinoid A, (4) and Frobenius give, in B=A\ensuremath{\otimes}E,

\[
       (\sum_j a_j\otimes e_j)^q=\sum_j a_j^q e_j^q\in A.    \tag{13}
\]

If f is a unit, both f\^{}q and (f\textsuperscript{\{-1\})}q lie in A
and multiply to one; so f\^{}q is an actual unit without reducing any
algebra. Any two analytic points above x have on every b the forced
value \textbar b\textsuperscript{q(x)\textbar{}}\{1/q\}. Scalar
surjectivity gives one, hence exactly one. Compact-to-Hausdorff
projection is therefore a homeomorphism on each affinoid; these
homeomorphisms glue on the ordinary affinoid interiors. Purity was
proved by prime chains in Section 3.2, without a reducedness assumption.

Let \ensuremath{\beta} be any ordinary (a,b) form near an upstairs point. The
homeomorphism and goodness give a downstairs affinoid neighborhood V
whose whole V\_E lies in its ordinary moment-description open. Its units
f\_j are elements of O(V)\ensuremath{\otimes}E. Put
g\_j=f\_j\textsuperscript{q\ensuremath{\in}O(V)}\emph{. Then t\_f=q\^{}\{-1\}t\_g\ensuremath{\circ}p\_E.
If \ensuremath{\omega} is the Euclidean coefficient superform, the downstairs form is
t\_g\^{}}(s\_q\^{}*\ensuremath{\omega}), where s\_q(t)=t/q. Its coefficient domain
contains t\_g(V) by scalar surjectivity. Explicitly, a basis term of
bidegree (a,b) is multiplied by q\^{}\{-(a+b)\}, as well as evaluating
its coefficient at t/q. Thus both prime and double-prime differential
factors descend. Equality as smooth sections follows either from this
moment identity or from compact-image nullity. The equality is on
int\_X(V) and its entire inverse image, including any zero region of \ensuremath{\beta}.

For a compact \ensuremath{\beta}, cover its compact image by finitely many such opens,
take the compact partition of Section 3.4 downstairs, multiply each
local descendant by its cutoff and extend by zero. The result \ensuremath{\alpha} is
compact, and its pullback is \ensuremath{\beta} globally: the cutoff sum is one near supp
\ensuremath{\beta} and \ensuremath{\beta} is zero elsewhere. No overlap agreement is needed for this
finite-sum construction. For completeness pullback is locally faithful:
in any downstairs unit chart and compact V, surjectivity gives identical
compact moment images for V and V\_E, so 3.4.4 says zero pullback
implies zero section on V. Hence local descendants also agree if sheaf
gluing is used. Under the homeomorphism their supports correspond. This
proves exact purely inseparable descent for every bidegree, with compact
supports. Applying (2) to D\ensuremath{\alpha} proves the finite purely inseparable
current equivalence against all complementary tests. No curve comparison
is involved.

\subsubsection{6.3. The actual embedded mixed
tower}\label{the-actual-embedded-mixed-tower}

For finite algebraic E/K in characteristic p, choose generators a\_i.
Write each irreducible minimal polynomial
f\_i(T)=g\_i(T\textsuperscript{\{p}\{r\_i\}\}), removing powers of p
until g\_i'\ensuremath{\ne}0. Substitution proves g\_i irreducible; coprimality with
its derivative proves separability. Thus
b\_i=a\_i\textsuperscript{\{p}\{r\_i\}\} are separable over K. Their
compositum S=K(b\_i)\ensuremath{\subseteq}E is finite separable: at each adjoining step its
minimal polynomial divides a separable polynomial, and successively
extending embeddings gives degree-many embeddings of the tower. Every
intermediate simple field is separable by comparing its number of
embeddings with the tower degree. For r=max r\_i, taking p\^{}r powers
of a rational expression gives E\textsuperscript{\{p}r\}\ensuremath{\subseteq}S. Thus the
\textbf{embedded} tower K\ensuremath{\subseteq}S\ensuremath{\subseteq}E is separable then purely inseparable. In
characteristic zero the derivative test makes every finite stage
separable. Finite tensor associativity with the bounded basis maps in
(4) identifies the actual spaces (X\_S)\_E=X\_E. Their componentwise
purity was proved in Section 3.2. Applying Sections 6.1 and 6.2 proves
zero-current transport to every finite embedded E\ensuremath{\subseteq}M. This uses no
previous inseparable or curve theorem as a premise.

\subsection{7. Completion, exact rational domains, and ordinary
presentation
opens}\label{completion-exact-rational-domains-and-ordinary-presentation-opens}

\subsubsection{7.1. Density and finite associativity in the actual
Banach
algebra}\label{density-and-finite-associativity-in-the-actual-banach-algebra}

M is the directed union of its finite embedded subextensions. For A\_L=A
completed-tensor\_K L, the \textbf{image} of A\ensuremath{\otimes}M is dense. First
approximate by finitely many algebraic tensors \ensuremath{\sum}a\_j\ensuremath{\otimes}\ensuremath{\lambda}\_j, \ensuremath{\lambda}\_j\ensuremath{\in}L; then
replace each \ensuremath{\lambda}\_j by m\_j\ensuremath{\in}M. The Banach tensor error is at most
max\_j\textbar\textbar a\_j\textbar\textbar{}
\textbar \ensuremath{\lambda}\_j\ensuremath{-}m\_j\textbar. This concerns actual Banach norms and
images; no injection of an infinite uncompleted tensor is presumed.

For finite E/K the maps (a\ensuremath{\otimes}e)\ensuremath{\otimes}\ensuremath{\lambda}\ensuremath{\mapsto}a\ensuremath{\otimes}e\ensuremath{\lambda} and a\ensuremath{\otimes}\ensuremath{\lambda}\ensuremath{\mapsto}(a\ensuremath{\otimes}1)\ensuremath{\otimes}\ensuremath{\lambda} extend to inverse
bounded maps

\[
        (A\widehat\otimes_K E)\widehat\otimes_E L
                          \simeq A\widehat\otimes_K L.      \tag{14}
\]

The first is bounded using the finitely many bounded coordinate maps of
E and the tensor bound; the second is bounded by the inclusion A\ensuremath{\to}A\_E
and the tensor bound. They are inverse on dense elementary tensors,
hence inverse on the completions. This applies also to any actual
rational localization affinoid, with arbitrary real radii. Density
supplies approximants, while the next argument, not density alone,
supplies exact domains and presentation opens.

Scalar surjectivity here can also be reconstructed. Given x with H=H(x),
embed the algebraic M into an algebraic closure of H over K, by a
maximal partial embedding and adjoining a root if it is not yet defined
on M. The finite extensions of H have their compatible unique absolute
values by Section 3.1. Complete that algebraic closure to a field F. The
restriction on each finite K-subfield of M is its specified norm, so the
embedding extends isometrically to L\ensuremath{\to}F. Evaluation of any tensor
representation obeys
\textbar\ensuremath{\sum}a\_j(x)\ensuremath{\lambda}\_j\textbar\ensuremath{\le}max\_j\textbar\textbar a\_j\textbar\textbar{}
\textbar \ensuremath{\lambda}\_j\textbar. Taking the infimum over representations bounds it
by the full tensor seminorm; it extends to a multiplicative seminorm of
A\_L over x, including the separated quotient in its completion. The
same proof over E works on every scalar-extended affinoid domain. This
avoids an assumed isometry or injectivity of A\ensuremath{\otimes}M. Compact inverse images
were proved in Section 3.3.

\subsubsection{7.2. Rational inequalities, equality loci, and strict
loci}\label{rational-inequalities-equality-loci-and-strict-loci}

Let A be affinoid over a finite E\_0\ensuremath{\subseteq}M, B=A\_L, and let R\_L\ensuremath{\subseteq}M(B) be the
rational domain \textbar f\_i\textbar\ensuremath{\le}r\_i\textbar g\textbar,
r\_i\textgreater0, with no common analytic zero. Put

\[
 \delta=\min_{y\in\mathcal{M}(B)}
    \max(|g(y)|,r_1^{-1}|f_1(y)|,\ldots,r_s^{-1}|f_s(y)|)>0. \tag{15}
\]

Compactness and no common zero give positivity. Choose approximants
g\_0,f\_\{i,0\} from one finite A\_E whose errors in B satisfy
\textbar\textbar g\_0\ensuremath{-}g\textbar\textbar\textless \ensuremath{\delta} and
\textbar\textbar f\_\{i,0\}\ensuremath{-}f\_i\textbar\textbar\textless r\_i\ensuremath{\delta}. Density
and (14) give this simultaneously. If \textbar g(y)\textbar\ensuremath{\ge}\ensuremath{\delta}, its norm
is unchanged. A numerator at or above r\_i\textbar g(y)\textbar{} also
has unchanged norm; a numerator below that threshold remains strictly
below, by ultrametricity. Thus every weak, equal, and strict comparison
is preserved at all points where \textbar g\textbar\ensuremath{\ge}\ensuremath{\delta}. If
\textbar g\textbar\textless \ensuremath{\delta}, some numerator has weighted norm at least
\ensuremath{\delta}. Its norm is unchanged, while \textbar g\_0\textbar\textless \ensuremath{\delta}, so
that inequality remains strictly false. In particular both the entire
weak system and the entire strict system are preserved on all M(B).
\textbf{Every equality locus inside the validity domain} is preserved;
no claim is made about irrelevant individual comparisons at points where
the whole system is false. Zeros and nilpotent elements cause no
exception, since approximation used the actual Banach norm.

The approximant weighted maximum is positive upstairs. Scalar
surjectivity implies no common analytic zero downstairs at E; spectral
surjectivity then implies they generate the unit ideal. They define a
genuine rational domain R\_E. Its base change has precisely the subset
R\_L. The universal analytic-domain property identifies their analytic
structures, because each domain inclusion factors uniquely through the
other. Thus R\_L=(R\_E)\_L as actual affinoid domains, including
nilpotents. The same argument gives exact inverse image equality for the
strict loci

\[
 O_E=\{|f_{i,0}|<r_i|g_0|\ \forall i\},\qquad O_L=q_E^{-1}O_E. \tag{16}
\]

O\_E is an ordinary open in the ambient affinoid and lies inside R\_E.
It is an ordinary interior presentation neighborhood. We do not assert
equality of all ordinary interiors or all absolute analytic interiors
under base change. Only these exact saturated strict loci are needed. A
vacuous system is handled by the harmless strict comparison
\textbar1\textbar\textless2\textbar1\textbar.

For each y and each prescribed ordinary description open upstairs, the
primary rational-neighborhood statement supplies an R\_L inside it with
y in its strict locus. Equations (15)--(16) produce a finite-stage
saturated presentation neighborhood of y. This establishes the needed
topological refinement directly; no inverse-limit theorem is imported.

\subsection{8. Exact descent and compact gluing of all higher-form
tests}\label{exact-descent-and-compact-gluing-of-all-higher-form-tests}

Let \ensuremath{\beta} be an ordinary (a,b) form near y\ensuremath{\in}X\_L, and x its image in X.
Choose a downstairs strict affinoid neighborhood V\_0 and an ordinary
U\_0\ensuremath{\subseteq}int\_X(V\_0) containing x. Intersect its upstairs inverse image
with a moment-description open for \ensuremath{\beta}. Shrink so that on this
intersection

\[
 \beta=\sum_{|I|=a,|J|=b}
       \omega_{IJ}(t_f)\,d't_{f,I}\wedge d''t_{f,J},          \tag{17}
\]

where all f\_j are actual analytic units and all coefficients are smooth
on a Euclidean open containing the moment image of the retained compact
neighborhood. Ambient coefficients exist by the stated closed-subset
extension theorem. Select a rational R\_L of y contained in this open
with y in its strict locus. Section 7 gives R\_L=(R\_\{E\_0\})\emph{L
and its descended strict locus. Since R\_L\ensuremath{\subseteq}U}\{0,L\}, scalar
surjectivity on the domain implies R\_\{E\_0\}\ensuremath{\subseteq}U\_\{0,E\_0\}. Its strict
locus is therefore an ordinary open of the \textbf{whole} X\_\{E\_0\},
inside its affinoid interior, rather than only a relatively open subset
at an exposed boundary.

On R\_L put \ensuremath{\delta}\_j=min\textbar f\_j\textbar\textgreater0. In its actual
affinoid Banach algebra, approximate all f\_j by u\_j from one finite
R\_E, E\ensuremath{\supseteq}E\_0, with
\textbar\textbar u\_j\ensuremath{-}f\_j\textbar\textbar\textless \ensuremath{\delta}\_j. Equations (14)
and density apply to this rational localization, not just to the
original A. Ultrametricity gives exactly

\[
                      |u_j|=|f_j|\quad\text{on }R_L.        \tag{18}
\]

Scalar surjectivity implies no zero of u\_j on R\_E, so u\_j is an
actual unit by the prime-spectrum criterion. This proof ensures its
actual inverse, without requiring spectral smallness of nilpotents.
Surjectivity also gives equality of compact moment images
t\_u(R\_E)=t\_f(R\_L), so the same coefficient open and \ensuremath{\omega}\_\{IJ\} are
valid at the finite stage.

Equation (18) is sufficient for higher forms for a precise reason. On
the upstairs ordinary presentation open, the unit f\_j/u\_j has
logarithmic scalar section identically zero. A numerical-zero smooth
scalar section is zero in the smooth sheaf: in each compact moment
description its coefficient vanishes on the moment image, and 3.4.4
gives nullity, followed by ordinary local sheaf gluing. Therefore
t\_\{f\_j\}=q\_E\^{}\emph{t\_\{u\_j\} \textbf{as smooth sections}, and
applying each derivation gives d't\_\{f\_j\}=q\_E\^{}}d't\_\{u\_j\} and
d'\,'t\_\{f\_j\}=q\_E\^{}\emph{d'\,'t\_\{u\_j\}. This also follows in
the joint (f,u) moment whose image lies on the diagonal. Substitution in
(17), including every exterior factor, proves exact equality
\ensuremath{\beta}=q\_E\^{}}b\_E on O\_L=q\_E\^{}\{-1\}O\_E, with b\_E the ordinary (a,b)
stage form obtained by replacing every f\_j by u\_j. This is not C\^{}0
convergence of differentiated forms. The equality holds on the entire
saturated open, including the regions where \ensuremath{\beta} is zero.

Now take arbitrary compact \ensuremath{\beta} on X\_L and F=supp \ensuremath{\beta}. If F is empty take
zero at K. Otherwise finitely many opens of this kind cover F. Enlarge
their finite fields to one E\ensuremath{\subseteq}M, and use (14) to base change all domains,
strict loci, and local forms. Thus O\_\{i,L\}=q\_E\^{}\{-1\}O\_\{i,E\}
and \ensuremath{\beta}=q\_E\^{}*b\_\{i,E\} on each entire such open. The compact
F\_E=q\_E(F) is covered by O\_\{i,E\}. Choose finite compact \ensuremath{\theta}\_i
supported in O\_\{i,E\} with sum one near F\_E by Section 3.4. Define

\[
             \alpha_E=\sum_i\widetilde{\theta_i b_{i,E}},    \tag{19}
\]

where tilde is extension by zero. Each product has support contained in
the compact closed set supp \ensuremath{\theta}\_i\ensuremath{\subseteq}O\_\{i,E\}. At every exterior point it
has an ordinary zero neighborhood, so this is a smooth compact form,
with no boundary derivative. Pullback of its summand is q\_E\^{}*\ensuremath{\theta}\_i \ensuremath{\beta}
on O\_\{i,L\}; off that open both sides are zero. Consequently,
globally,

\[
           q_E^*\alpha_E=(\sum_i q_E^*\theta_i)\beta=\beta.  \tag{20}
\]

The sum is one near F and \ensuremath{\beta} is zero outside F. The support of \ensuremath{\alpha}\_E is
contained in a finite union of compact cutoff supports; their full
scalar inverse image is compact by Section 3.3. We have neither assumed
the original support saturated nor silently glued higher forms from a
scalar lemma. Local descriptions agree with \ensuremath{\beta} on their full preimage
opens, and (19) is one concrete smooth global form. Agreement of b\_i on
stage overlaps is unnecessary for this construction. Differentiation of
the exact section identity (20) proves Theorem A in every bidegree.

One can separately check local faithfulness if desired. In a downstairs
unit chart and compact affinoid V, scalar surjectivity gives identical
moment images for V and V\_L. Applying compact-image nullity upstairs
and downstairs shows that zero pullback implies zero section on V.
Ordinary neighborhoods then give faithfulness on all ordinary opens.
This is a proved consequence of compact-image equality, not blanket
invariant-form descent, and is not needed for the finite-sum gluing.

\subsection{9. Independent componentwise purity at
completion}\label{independent-componentwise-purity-at-completion}

We do not use simultaneous Noether normalization, flat completed
tensoring of an injection, or a claim inferred from total dimension
alone. Instead, exact rational-domain descent makes the following
contradiction available in every dimension.

Fix y\ensuremath{\in}X\_L and a strict affinoid neighborhood V\_0 of its downstairs
image x, with y\ensuremath{\in}(int\_X(V\_0))\_L. Its scalar extension has total
dimension n by the primary dimension theorem, so dim\_y(X\_L)\ensuremath{\le}n.~Suppose
its local dimension is m\textless n.~In (V\_0)\_L, remove the finitely
many minimal components of dimension greater than m. None contains y by
the primary local dimension formula. The remaining ordinary open
contains y. Choose a rational neighborhood R\_L of y inside it and
inside (int\_X(V\_0))\_L, with strict inequalities at y. Every local
dimension of this domain is at most m, again by the component/domain
formula; hence dim(R\_L)\ensuremath{\le}m.

Section 7 descends R\_L exactly to a nonempty R\_E in (V\_0)\_E at one
finite E\ensuremath{\subseteq}M. Scalar surjectivity on R\_E places it in (int\_X(V\_0))\_E,
hence in the actual X\_E. That space is pure n by Sections 3 and 6, so
the domain formula gives dim(R\_E)=n, even if the auxiliary rational
radii make this domain non-strict. Complete scalar dimension invariance
now gives

\[
                    \dim R_L=\dim(R_E)_L=\dim R_E=n,
\]

contradicting dim R\_L\ensuremath{\le}m\textless n.~Therefore dim\_y(X\_L)=n at every
point. This argument controls local dimensions, not just the total
dimension. For every minimal prime Q of any strict upstairs affinoid
domain, spectral surjectivity gives a point with kernel Q on that
component alone. The local component formula then yields dim(A\_L/Q)=n.
Thus \textbf{every} actual upstairs minimal component has dimension n.

Nilpotents were never removed from the rational algebras or their
completed tensors. The only reductions earlier were finite separable
prime calculations for dimension. The completed-stage purity proof does
not assert completed scalar extension preserves arbitrary module
injections. Goodness, strictness, separatedness, boundarylessness, and
compact scalar inverse images were checked in Section 3.3 using their
actual presentations and primary boundary theorem. We may therefore use
the canonical level-n tropical structure on all finite and completed
stages, and its domain/top-measure compatibility.

\subsection{10. Pairing, lengths, and conclusion of the
gate}\label{pairing-lengths-and-conclusion-of-the-gate}

Assume B\_h=0. Sections 4--6 prove B\_\{h\_E\}=0 on every finite
embedded E\ensuremath{\subseteq}M, against all ordinary complementary forms. Given
\textbf{any} compact ordinary complementary \ensuremath{\beta} on X\_L, Theorem A
supplies \ensuremath{\beta}=q\_E\^{}\emph{\ensuremath{\alpha}\_E at one such E. Its derivative is
q\_E\^{}}D\ensuremath{\alpha}\_E, a smooth top form. Both fields' actual spaces are pure
n, and all supports are compact. Applying (2), then the integration
formula for signed Radon pushforward, gives

\[
 \begin{aligned}
 B_{h_L}(\beta)
 &=\int_{X_L}(h_E\circ q_E)(q_E^*D\alpha_E)_{\mathrm{Rad},L}\\
 &=\int_{X_E}h_E(D\alpha_E)_{\mathrm{Rad},E}
 =B_{h_E}(\alpha_E)=0.
 \end{aligned}                                             \tag{21}
\]

The continuous multiplier is bounded on the common compact. No
continuous-current pullback has been defined. This proves ascent. For
the converse, pull back any compact downstairs complementary \ensuremath{\alpha}; its
support remains compact, and (2) gives B\_\{h\_L\}(\ensuremath{\pi}\^{}*\ensuremath{\alpha})=B\_h(\ensuremath{\alpha}).
Vanishing upstairs implies vanishing downstairs, proving (T).

The canonical measure formula preserves all actual Artinian lengths. If
R is any finite Artinian C-algebra of vector-space dimension d and C'/C
any complete extension, its algebraic tensor is finite-dimensional,
already complete, of C'-dimension d.~Decomposing it into local factors
R\_j and taking a composition series in each gives

\[
 d=\sum_j\dim_{C'}R_j
   =\sum_j\operatorname{length}_{R_j}(R_j)[\kappa_j:C'].      \tag{22}
\]

There is no separability or reducedness assumption. A finite-flat moment
fibre is this actual algebra. CLD 13.3.3 and the full primary 14.3.6
proof choose the canonical lift of the target torus Gauss point and sum
precisely these local degrees on the skeletal n-cells. Equation (22)
preserves that sum; the positive calibrations then give (2). Pointwise
generic multiplicities may change, and are not divided out. In contrast
same-carrier field restriction composes the moment with a finite-flat
torus projection of rank {[}F:K{]}, giving (3), used only in (6). These
two comparisons remain different in wild and inseparable cases.

The complete claim proved in this draft is the new completed-algebraic
\textbf{current transport gate}, with exact higher-test descent and
independent purity. The finite-Galois theorem is supporting mathematics,
not reported as a new theorem by repetition. Relative to the supplied
catalogue, the gate applies to every actual pure singular strict germ,
without its coordinate-product, invariant, real-PL, factorization, or
imposed upstairs-current hypotheses. Nevertheless it gives no
recognition or two-sign approximation in these germs. Catalogue
descriptions are not being used to infer those conclusions. Arbitrary
transcendental-extension transport, arbitrary analytic curve pullback,
and ordinary two-sign regularization remain separate unproved
obligations. The original non-strict arbitrary-radius negative example
is not altered and is not an original-scope counterexample. The status
for original P06.2 is partial, and candidate=false.

\subsection{11. Review, source boundaries and reproducibility}

The complete new mathematical core, both canonical native notes, author
audit and research limits were read before publication. All thirty-three
primary bodies were checked byte for byte against authentic origins and
the identical previously fully read bodies; this was exact source reuse,
rather than a new fresh reading of the entire quotation appendix.
Every new application was assessed in the present proof. The independent
Xean native Codex review used \texttt{gpt-6.1-sol} with \texttt{xhigh}
reasoning and returned PASS for Theorems A and B. The complete report,
all twenty-three premises and all thirty-two passages were read and
their exact indices, statements and immutable proof bindings checked.

The complete quoted primary units and their printed proofs remain in
the accompanying frozen argument. Source-internal references remain
normal primary references; this paper imports the exact stated analytic
foundations, rather than claiming to reprove their entire underlying
theories. The reviewer verified the essential CLD and published Ducros
sources independently. CLD's composition statement supplies the needed
boundary facts; the additional Berkovich citation adds no indispensable
premise. Base change preserves interior points; no unrestricted equality
of relative interiors is used.

Harmless source-display corrections, already avoided by this proof,
are recorded separately: graded commutativity uses
\((p+q)(p'+q')\); rational localization retains its ambient algebra;
relative-interior composition uses \(g^{-1}(\operatorname{Int}(Y/X))\);
the Artinian case has \(\dim_y(Y)=0\); the target volume in scalar
projection is downstairs volume; CLD 5.10.1(a)'s comparison is equality;
and the cutoff has the correct values near the two function images,
with its neighborhood inside a point-separating open. Section 3.4 supplies
that cutoff proof independently. Frozen raw quotations are unchanged.

The small rational-norm, exterior-basis and unreduced-algebra diagnostics
are finite reproducibility evidence. They do not prove the analytic
theorems. The standalone source was submitted to the native LaTeX editor
for compilation. Same-model review, computation and compilation are
not formal verification or priority certification. Original classical
recognition and the further bridges explicitly excluded above remain open.

\subsection{Primary sources and accompanying argument}

Chambert-Loir and Ducros, \emph{Formes diff\'erentielles r\'eelles et
courants sur les espaces de Berkovich},
\url{https://arxiv.org/abs/1204.6277v2}; Ducros,
\emph{Espaces de Berkovich, polytopes, squelettes et th\'eorie des
mod\`eles}, Confluentes Math. 4 (2012), 1250007,
\url{https://www.numdam.org/item/10.1142/S1793744212500077/};
Berkovich, \emph{\'Etale cohomology for non-Archimedean analytic spaces},
Publ. Math. IHES 78 (1993), 5--161.

The accompanying frozen argument retains all thirty-three full selected
primary bodies. Every load-bearing supporting mathematical proof and
the exact imported foundations above are retained in this rendering.

\end{document}
